\subsubsection{(iii) Gravitational Radiation from the Quasi-normal Modes}

The power radiated in gravitational waves by the quasi-normal pulsations is equal
to the time rate of change of the fluid's pulsation energy. (Here we refer to energies
measured at infinity with the energy redshift effect taken into account; cf.\ Harrison,
Thorne, Wakano, and Wheeler [1965], pp.\ 19--21; also RSSD \S~3.1.2.) It is easy to
write down an expression for the kinetic energy of fluid motion but very difficult to
construct an expression for the compressional and gravitational potential energy.
Therefore, we shall calculate only the kinetic energy, and we shall then argue that the
potential energy must be equal to the kinetic energy but must pulsate out of phase with
it.

\subsubsection{(iv) Pulsation Energy and Power Radiated}

The kinetic energy of pulsation is given by (cf.\ RSSD, eq.\ [4.26'])
%
\begin{equation}
E_{\rm kin} = \int_{\rm star}
\underbrace{\left(\rho + p\right)}_{\text{inertial mass}\atop\text{per unit volume}}
\underbrace{\left(\frac{1}{2}v^2\right)}_{\text{physical velocity}\atop\text{of fluid}}
\underbrace{e^{(\nu - \lambda)/2}}_{\text{gravitational}\atop\text{redshift}}
\underbrace{d\mathcal{V}}_{\text{proper}\atop\text{volume}}
= \int_0^R \frac{1}{4}(\rho+p)\,e^{(\lambda-\nu)/2}\bigl[r^{-2}\xi_r^2(r) + (r^{-1}\xi_\perp(r))^2\bigr]\,e^{\nu/2}\,4\pi r^2\,dr\,d\theta\,d\phi \,.
\tag{25}
\end{equation}
%
Upon substituting equations (7a), (11), (12), and (22) into this expression and
integrating over $\theta$, we obtain for the fluid kinetic energy of the $n$th quasi-normal mode
%
\begin{equation}
E_{\rm kin} = 2\pi\,(2l+1)^{-1}\,\sigma_n^2\,e^{-2i\sigma_n t}
\int_0^R (\rho+p)\,e^{(\lambda-\nu)/2}
\Bigl\{
r^{-2}\bigl|W_n^{(l)}(r)\bigr|^2
+ l(l+1)\,r^{-2}\bigl|V_n^{(l)}(r)\bigr|^2
\Bigr\}
\Bigl[
1 + l(l+1)\,\bigl|V_n^{(l)}(r)\bigr|^2
\sin^2\bigl(\sigma_n t + \delta_n(r)\bigr)
+ \bigl\{l + \delta_n(r)\bigr\}
\Bigr]\,dr \,,
\tag{26}
\end{equation}
%
Here $\delta_n(r)$ and $\delta_n(r)$ are the phases of $W_n^{(l)}(r)$ and $V_n^{(l)}(r)$, and we have assumed that
$1/\tau_n \ll \sigma_n$.

For situations of physical interest---e.g., neutron stars formed in a supernova explosion or neutron stars undergoing relaxation oscillations---rough estimates suggest
that $1/\tau_n \ll \sigma_n$ (see Wheeler 1966; Zee and Wheeler 1967; Misner and Zapolsky 1967;
Chau 1967). When this is true, the eigenfunctions (14) couple the real and imaginary
parts of the eigenfunctions only very slightly; and, consequently, one has
%
\begin{equation}
\delta_n(r) \ll 1, \quad \delta_n(r) \ll 1 \,.
\tag{27}
\end{equation}
%
and
%
\begin{equation}
E_{\rm kin} = E_{\rm pulse}^n \sin^2(\sigma_n t) \,.
\tag{28a}
\end{equation}
%
Here
%
\begin{equation}
E_{\rm pulse}^n = 2\pi\,(2l+1)^{-1}\,\sigma_n^2
\int_0^R (\rho+p)\,e^{(\lambda-\nu)/2}
\Bigl[
r^{-2}\bigl|W_n^{(l)}(r)\bigr|^2
+ l(l+1)\,r^{-2}\bigl|V_n^{(l)}(r)\bigr|^2
\Bigr]\,dr \,.
\tag{28b}
\end{equation}
%
When $1/\tau_n \ll \sigma_n$, the total fluid pulsation energy---kinetic plus potential---is very
nearly constant over a total pulsation period, $\Delta t = 2\pi/\sigma_n$. This is possible only if the
potential energy of pulsation is given by
%
\begin{equation}
E_{\rm pot} \approx E_{\rm pulse}^n \cos^2(\sigma_n t) \,,
\tag{29a}
\end{equation}
%
so that
%
\begin{equation}
E_{\rm kin} + E_{\rm pot} = E_{\rm pulse}^n \,.
\tag{29b}
\end{equation}
%
is nearly time independent. Equations (29a) and (29b) for the total pulsation energy
should be roughly correct even when $1/\tau_n$ is not small compared to $\sigma_n$.
